\documentclass[11pt,a4paper]{article}
\usepackage[margin=25mm,headheight=16pt]{geometry}
\usepackage{fontspec,kotex,xcolor,hyperref,fancyhdr}
\setmainfont{DejaVu Serif}
\setmainhangulfont{Noto Sans CJK KR}
\definecolor{navy}{HTML}{173E59}
\pagestyle{fancy}\fancyhf{}
\fancyhead[L]{\small GuardSynth CoC}
\fancyhead[R]{\small 제목·초록·서론 · 한글}
\fancyfoot[L]{\small v0.1 · 2026-09-30}
\fancyfoot[R]{\small\thepage}
\renewcommand{\headrulewidth}{0.3pt}
\setlength{\parindent}{0pt}
\setlength{\emergencystretch}{3em}
\linespread{1.15}\clubpenalty=10000\widowpenalty=10000
\hypersetup{pdftitle={GuardSynth: CoC 기반 VLM 학습 보강을 위한 행동 제약의 정형 명세와 검증},pdfsubject={Title and abstract for author review},pdfauthor={GuardSynth project}}
\hypersetup{colorlinks=true,urlcolor=navy}\begin{document}
\vspace*{8mm}
{\Large\bfseries\color{navy}\raggedright GuardSynth: CoC 기반 VLM 학습 보강을 위한 행동 제약의 정형 명세와 검증\par}
\vspace{8mm}
{\large\bfseries 초록\par}
\vspace{3mm}
Alpamayo와 같은 추론 기반 자율주행 모델은 Chain of Causation(CoC)을 통해 장면에 대한 인과적 추론과 행동 예측을 연결한다. 그러나 상황과 행동의 이유를 설명하는 CoC 기반 학습만으로는 올바른 행동 선택과 제약 준수를 충분히 학습하는 데 한계가 있다. 본 논문은 이러한 한계를 보완하기 위해 행동 제약을 체계적으로 명세·검증하고 CoC 학습 데이터에 결합하는 GuardSynth를 제안한다. 행동 제약의 범위·특성·구성 요소를 정의하고, 적용 조건, 금지 행동, 유지·해제 조건, 시간 및 근거를 Evidence-Bound Lifecycle Contracts(EBLC)로 명세한다. 이어 제한된 논리 모델에서 일관성과 지정 속성을 검증하고, 명세에 대응하는 통제 자연어(Controlled Natural Language, CNL)를 생성하여 CoC를 보강한다. 통제된 합성 시각·시간 과제에서 시각언어모델(VLM)인 Qwen3-VL-2B-Instruct를 동일한 LoRA 학습 예산과 세 개의 시드로 비교하였으며, 제약 추가 조건에는 학습과 평가 모두에서 제약을 제공하였다. CoC 단독, 기존 자연어 제약 추가, EBLC 기반 CNL 추가 조건의 평균 행동 선택 정확도는 각각 49.86\%, 96.81\%, 98.06\%였으며, 제약 위반율은 19.17\%, 1.39\%, 0.83\%였다. 결과는 행동 제약의 정형 명세·검증을 CoC 학습에 연결하는 체계와 명시적 제약 제공에 따른 행동 선택 개선을 보여준다.
\clearpage
\section*{I. 서론}
자율주행의 행동 판단에는 장면에 존재하는 객체를 인식하는 것과 함께, 객체의 관계와 상황의 변화를 고려하여 적절한 행동을 선택하는 과정이 필요하다. NVIDIA의 Alpamayo-R1은 이러한 판단을 위해 Chain of Causation(CoC) 추론과 주행 궤적 예측을 결합하는 시각·언어·행동 모델(VLA)을 제시하였다. 인과적으로 연결된 설명을 운전 행동에 대응시키고, 지도 미세조정과 강화학습을 통해 추론 품질과 추론–행동 간 일관성을 높이는 접근이다. 이는 운전 행동뿐 아니라 그 행동에 이르는 판단 과정도 학습에 활용하는 방향을 보여준다.~\href{\detokenize{https://arxiv.org/abs/2511.00088}}{[1]}\par\vspace{5pt}
관련 연구에서도 시각 정보와 언어적 추론을 연결하는 다양한 방법이 제안되었다. DriveLM은 지각·예측·계획에 관한 질문과 답변을 그래프로 연결하여 운전 상황의 다단계 추론을 구성하였다. DriveVLM은 장면 설명, 장면 분석, 계층적 계획을 결합하여 복잡한 운전 상황을 처리하였다. 이들 연구는 장면 해석을 행동 판단과 계획에 활용하는 기반을 제공한다. 본 연구는 이 흐름에서 학습 데이터에 제공하는 행동 제약의 표현과 검증에 초점을 둔다.~\href{\detokenize{https://arxiv.org/abs/2312.14150}}{[2]},~\href{\detokenize{https://arxiv.org/abs/2402.12289}}{[3]}\par\vspace{5pt}
그러나 상황과 행동의 이유를 설명하는 CoC 기반 학습만으로는 올바른 행동 선택과 제약 준수를 충분히 학습하는 데 한계가 있다. 행동의 이유를 설명하는 정보와 행동의 허용 범위를 규정하는 정보는 역할이 다르기 때문이다. 예를 들어 “보행자에게 양보하기 위해 감속한다”는 설명은 행동의 이유를 전달하지만, “횡단 영역이 점유된 동안에는 진입하지 않는다”는 제약은 선택 가능한 행동을 제한한다. 영역이 비워진 뒤 요구되는 대기 시간이 달라지면, 같은 장면에서도 허용되는 진입 시점이 달라질 수 있다. 본 연구에서 다루는 문제는 이러한 상황 의존적 제약을 설명 중심의 학습 데이터에 명시적으로 결합하여 행동 선택을 보강하는 것이다.\par\vspace{5pt}
우리의 선행 자연어 제약 실험에서는 통제된 합성 시간 과제에서 CoC형 설명에 제약을 추가하였을 때 조기 진입을 줄이고 제약을 준수하는 목표 완료를 개선하는 결과를 관찰하였다. 해당 과제에는 동일한 장면에 서로 다른 대기 조건이 적용되는 경우가 포함되므로, 제약을 제공하는 것은 행동 선택에 필요한 정보를 추가하는 의미를 갖는다. 이 결과를 출발점으로, 본 연구는 제약을 단순히 문장으로 추가하는 것에서 나아가 제약의 구성과 의미를 체계적으로 다룬다.~\href{\detokenize{../../../../../../projects/01-safety-constrained-coc/docs/reports/TEMPORAL_GUARD_MULTISEED_EXPERIMENT_REPORT_V01.md}}{[4]}\par\vspace{5pt}
이를 위해 세 가지 질문에 답해야 한다. 첫째, CoC를 보강하는 행동 제약의 범위와 구성 요소는 무엇인가? 둘째, 제약이 표현하는 조건과 행동 제한의 일관성을 어떤 논리 모델에서 확인할 것인가? 셋째, 확인한 명세를 학습에 사용할 자연어로 어떻게 전달할 것인가? 자유로운 자연어 표현은 사람이 이해하기에 유용하지만, 자동 검증에 필요한 의미와 가정을 명시하려면 별도의 구조가 필요하다. 따라서 본 연구는 행동 제약의 명시성·추적성·논리적 일관성과 자연어 학습 입력의 연결을 함께 다루는 체계를 제안한다.\par\vspace{5pt}
제안하는 GuardSynth는 행동 제약의 범위·특성·구성 요소를 정의하고, 주체와 대상, 적용 조건, 금지 행동, 유지·해제 조건, 시간 및 근거를 Evidence-Bound Lifecycle Contracts(EBLC)로 명세한다. 선택한 프로파일을 Core/SMT 표현으로 변환하여 제한된 모델에서 일관성과 지정 속성을 검증하고, 명세 필드에 대응하는 통제 자연어(Controlled Natural Language, CNL)를 생성하여 CoC에 결합한다. 정형 명세는 자연어를 대체하는 최종 학습 입력이 아니라, 학습에 제공할 제약을 구조화하고 그 의미를 확인하기 위한 중간 표현으로 사용한다. 명세 검증과 제약 보강의 행동 효과는 각각 논리적 근거와 경험적 근거로 구분하여 평가한다.\par\vspace{5pt}
행동 효과는 통제된 합성 시각·시간 과제에서 Qwen3-VL-2B-Instruct를 동일한 LoRA 학습 예산과 세 개의 시드로 학습하여 평가한다. CoC 단독(P0), 우리의 기존 자연어 제약 추가(P1), EBLC 기반 CNL 추가(P2)를 비교하며, P1과 P2에는 학습과 평가 모두에서 제약을 제공한다. 평균 행동 선택 정확도는 각각 49.86\%, 96.81\%, 98.06\%이고, 제약 위반율은 19.17\%, 1.39\%, 0.83\%이다. 주 비교는 P0 대비 제약 보강 조건의 효과이며, P1과 P2의 비교는 우리의 두 제약 표현 방식을 비교하는 내부 분석이다. 알파마요는 연구 배경이며, 본 논문의 학습·평가 대상은 Qwen 기반의 통제 과제이다.~\href{\detokenize{../temporal-experiment-closure-2026-09-29-001/all_primary_metrics.csv}}{[5]}\par\vspace{5pt}
본 논문의 기여는 다음 세 가지이다.\par\vspace{5pt}
\begin{enumerate}
\item CoC 행동 제약의 범위·특성·구성 요소를 정의하는 제약 모델을 제안한다.
\item EBLC 명세·논리 검증·CNL 생성·CoC 보강으로 이어지는 방법을 제안한다.
\item 동일 예산의 P0·P1·P2 비교를 통해 명시적 제약 제공에 따른 행동 선택 정확도와 제약 준수의 개선을 평가한다.
\end{enumerate}
이후 II장에서는 관련 연구를, III장에서는 행동 제약 모델을, IV장에서는 GuardSynth의 명세·검증·CNL 생성 방법을 설명한다. V장과 VI장에서는 실험 설정과 결과를 제시하고, VII장과 VIII장에서는 결과의 의미와 적용 범위를 논의한 뒤 결론을 정리한다.\par\vspace{5pt}
\section*{참고문헌 및 집필 근거}
\begin{itemize}
\item [1] NVIDIA, “Alpamayo-R1: Bridging Reasoning and Action Prediction for Generalizable Autonomous Driving in the Long Tail,” arXiv:2511.00088, 2025, v2 2026.
\item [2] C. Sima et al., “DriveLM: Driving with Graph Visual Question Answering,” ECCV, 2024, arXiv:2312.14150.
\item [3] X. Tian et al., “DriveVLM: The Convergence of Autonomous Driving and Large Vision-Language Models,” arXiv:2402.12289, 2024.
\item [4] 연구팀 선행 실험 기록: STOP–HOLD–RELEASE–GO 시간적 Guard 다중 Seed 실험 보고서 v1, 2026-08-04. 내부 실험 보고서이며 출판 논문의 서지정보를 대신하지 않는다.
\item [5] 현재 연구의 확장 실험 원자료: temporal-experiment-closure-2026-09-29-001, all\_\allowbreak{}primary\_\allowbreak{}metrics.csv, expansion/test 집계.
\end{itemize}
\end{document}
